\section{Conventions and the reduction theorems}\label{sec:reductions}

We first specify the homological statement used in induction and the
classical groups to which the known reductions apply. Unless coefficients
are specified otherwise, homology groups and chains have coefficients in $\Q$.

\subsection{Posets, chains, and elementary rank}
For a finite poset $P$, let $\Delta(P)$ be its order complex and write
$\widetilde H_*(P)$ for its reduced homology. We use the augmented chain complex:
there is a one-dimensional group in degree $-1$, generated by the empty
simplex, and the boundary of each vertex is that generator. Consequently
$\widetilde H_{-1}(\varnothing)=\Q$. For a nonempty finite complex,
nonvanishing in degree $-1$ is impossible.

A \emph{full flag} in an elementary abelian group $E$ of rank $r$ is a
chain
\[
E_1<E_2<\cdots<E_r=E,
\qquad \rk E_i=i.
\]
We also call such a flag saturated. A coefficient at a full flag always
means its coefficient as an oriented simplex, with the displayed
increasing order. Write
\[
m_p(G)=\max\{\rk E:E\leq G\text{ an elementary abelian }p\text{-subgroup}\},
\]
allowing $E=1$ in this maximum. The order complex of $\Ap(G)$ has
dimension $m_p(G)-1$, including the empty case. We say that $G$ has the
\emph{Quillen dimension property} at $p$ if
\begin{equation}\label{eq:QD}
\widetilde H_{m_p(G)-1}(\Ap(G))\ne0.
\end{equation}
A nonzero cycle in that degree cannot be a boundary, because the chain
group in the next degree is zero.

The \emph{homological implication} for $(G,p)$ is
\begin{equation}\label{eq:HQC}
\Op(G)=1\quad\Longrightarrow\quad
\widetilde H_*(\Ap(G))\ne0.
\end{equation}
When arguing by minimality we fix $p$ and minimize $|G|$ among
counterexamples to \eqref{eq:HQC}. Thus every group of smaller order
satisfies \eqref{eq:HQC}, with no restriction on its composition factors.
If $p\nmid|G|$, its elementary poset is empty and it is not a
counterexample under our convention.

We use two standard elementary facts about finite posets. Pointwise
comparable order-preserving maps induce homotopic maps on order
complexes. Also, deleting a vertex whose strict upper interval is
contractible preserves homotopy type: its link is the join of its lower
and upper intervals and is contractible. These facts also follow from
the standard fiber and gluing lemmas for posets; see
\citet{Quillen1978,PitermanSmith2025}.

\subsection{Components and automorphisms}
A component of a finite group is a subnormal quasisimple subgroup.
The components supplied by the reduction below are simple. Distinct
components commute; distinct conjugates of a simple component intersect
trivially and generate their direct product. We use these standard
component facts and the usual automorphism structure of finite
classical simple groups; see \citet{GLS1998}.

If $L$ is nonabelian simple, a subgroup containing $L$ and acting
faithfully on $L$ by conjugation is identified with a subgroup of
$\Aut(L)$. An \emph{elementary $p$-extension} of $L$ means a split
extension $L\rtimes B$ in which $B$ is elementary abelian and induces
only outer automorphisms; $B=1$ is allowed. In particular, an elementary
quotient alone is not being taken as a splitting hypothesis. In our
configurations, $H=LA$ with $A$ elementary and $H$ faithful on $L$; a
vector-space complement to $A\cap L$ in $A$ supplies the required
subgroup $B$.

Every such faithful extension $L\le H\le\Aut(L)$ has $O_p(H)=1$.
Indeed, for $Q=O_p(H)$ simplicity gives $Q\cap L=1$, and normality
gives $[Q,L]\le Q\cap L$. Hence $Q\le C_H(L)=1$.
Thus these extensions meet the trivial-core condition in the
conditional dimension-property convention of the cited reductions.

For the linear and unitary groups, the parameter $q$ is the order of
the base field: the natural unitary space is over $\F_{q^2}$. In the
orthogonal sections, the discriminant means the determinant square
class of a Gram matrix. Its relation to the Witt sign includes the
usual dimension-dependent factor. These two conventions will not be
interchanged.

\subsection{The established reductions}

\begin{theorem}[Reduction at two]\label{thm:reduction-two}
Let $G$ be a minimal-order counterexample to \eqref{eq:HQC} for $p=2$.
Then $G$ has a simple component $L$ in the following list, where the
field parameter $q$ has characteristic at least five:
\[
\begin{array}{ll}
\PSL_n(q),\ \PSU_n(q)& n\geq4,\\
\PSp_{2n}(q)& n\geq3,\\
\Omega_{2n+1}(q)& n\geq2,\\
\mathrm P\Omega_{2n}^{\pm}(q)& n\geq4.
\end{array}
\]
\end{theorem}

This is the part of \citet[Theorem~B]{Piterman2026} that we require.
The theorem there also supplies an extension failing the dimension
property; our argument at two will instead use maximal faithful
configurations and does not need that additional conclusion.

\begin{theorem}[Odd-prime reduction]\label{thm:reduction-odd}
Suppose that, for every odd prime $p$, every odd prime power $q$ with
$p\mid q+1$, and every simple $L=\PSU_n(q)$ with $n\geq5$, each
elementary $p$-extension of $L$ inside $\Aut(L)$ has the Quillen
dimension property. Then \eqref{eq:HQC} holds for every finite group
and every odd prime.
\end{theorem}

\begin{proof}[Derivation from the established reduction]
Fix an odd prime and take a minimal counterexample, if one exists.
The inductive hypothesis in \citet[Theorem~1.4]{PitermanSmith2025}
holds by minimality: its proper subgroups and proper quotients by
central subgroups of order prime to $p$ have smaller order.
The reduction to simple components in its published proof (p.~367)
uses \citet[Theorem~2.22(2) and Lemma~2.3]{PitermanSmith2025}.
Theorem~1.4 then reduces the
remaining assertion to the dimension property for the indicated
unitary extensions occurring in the group. Proving it for all such
extensions inside the automorphism group suffices.

The dimension-property list in
\citet[Theorem~2.16(2)]{PitermanSmith2025}, recalled from
\citet[Theorem~3.1]{AschbacherSmith1993}, has possible unitary exceptions in this
range only when
\[
n\geq q(q-1),
\qquad\text{or}\qquad
n\geq s(s-1),\quad s=q^{1/p},
\]
where the second possibility requires a field order $s$ and a
nontrivial field extension. Since $q$ is odd and $p\mid q+1$,
$q\geq5$. In the second case $q=s^p$ gives
$s\equiv s^p\equiv-1\pmod p$, so $s\geq5$ as well.
Both thresholds are at least twenty; in particular no missing
dimension $n\leq4$ needs a new argument. The rank-one coincidences
$\PSU_2(5)\cong A_5$ and $\PSU_2(9)\cong A_6$
\citep[Theorems~10.9 and~4.6(ii),(iv)]{Taylor1992} introduce no
exception for the respective odd primes dividing $q+1$ in the same
list. The proposed hypothesis therefore supplies every remaining
case required by the reduction.
\end{proof}

The rest of the paper proves the new assertions needed by these two
theorems. We first establish the chain and propagation methods; no
dimension property at two is assumed in those methods.

Table~\ref{tab:local-targets} records where these local tasks are
completed. At two, a supported coefficient is transferred to the
ambient group by the propagation lemma; it need not occur in the
largest chain degree of the component extension. The final exceptional
case also has a branch using an equivalent poset instead of propagation.

\begin{table}[ht]
\centering
\small
\begin{tabular}{@{}>{\raggedright\arraybackslash}p{.11\linewidth}>{\raggedright\arraybackslash}p{.60\linewidth}>{\raggedright\arraybackslash}p{.20\linewidth}@{}}
\toprule
Prime & Local task & Result \\
\midrule
Odd $p$ & Top-degree homology for the unitary extensions in
Theorem~\ref{thm:reduction-odd} &
Proposition~\ref{prop:unitary-odd} \\[3pt]
$2$ & Linear and unitary components of natural dimension at least five &
Proposition~\ref{prop:linear-two} \\[3pt]
$2$ & Symplectic components on the reduction list &
Proposition~\ref{prop:symplectic} \\[3pt]
$2$ & Orthogonal components, including the dimension-four linear and
unitary groups through their six-dimensional realizations, apart from
the case in the next row &
Proposition~\ref{prop:orthogonal} \\[3pt]
$2$ & The remaining square-determinant six-space over $\F_5$,
with simple group $\PSL_4(5)$ & Proposition~\ref{prop:psl45} \\
\bottomrule
\end{tabular}
\caption{The local tasks exported to the reductions. All prime-two
rows have defining characteristic at least five. The propositions in
those rows exclude the indicated components of a minimal
counterexample; they do not assert a general dimension property.}
\label{tab:local-targets}
\end{table}
